\documentclass[11pt,a4paper]{article}

%% PACHETE MATEMATICE
\usepackage{amsmath,amssymb,amsfonts}
\usepackage{amsthm}
\usepackage{mathpazo}
\usepackage{eulervm}
\usepackage{enumerate}
\usepackage{color}
\usepackage{graphicx}
\usepackage[all]{xy}
\usepackage{xcolor}
\usepackage{framed}
\usepackage[unicode]{hyperref}
\setcounter{MaxMatrixCols}{20}
\usepackage{multicol}
%\usepackage[romanian]{babel}


%% ASPECT
\linespread{1.05}
\usepackage[T1]{fontenc}
\usepackage[utf8x]{inputenc}
\usepackage[margin=2cm]{geometry}
% \usepackage[color]{showkeys}
\usepackage{anyfontsize}
\usepackage{titlesec}
\titleformat*{\section}{\large\bfseries}

\usepackage{fancyhdr}
\pagestyle{fancy}
\rhead{\small \color{gray} Notițe de Adrian Manea}
\lhead{\small \color{gray} AM-ETTI, Prof. C. Ghiu}


\usepackage[autostyle = false, style = german]{csquotes}
\usepackage[romanian]{babel}
\newcommand\qq{\enquote}

%% COMENZILE MELE
\renewcommand*{\proofname}{Dem.:}
\newcommand\bb{\mathbb}
\newcommand\dr{\mathrm}
\newcommand\kal{\mathcal}
\newcommand\fk{\mathfrak}
\newcommand\op{\oplus}
\newcommand\ot{\otimes}
\newcommand\ds{\displaystyle}
\newcommand\id{\indent}
\newcommand\nid{\noindent}
\newcommand\seq{\subseteq}
\newcommand\sneq{\subsetneq}
\newcommand\speq{\supseteq}
\newcommand\spneq{\supsetneq}
\newcommand\rar{\rightarrow}
\newcommand\Rar{\Rightarrow}
\newcommand\xrar{\xrightarrow}
\newcommand\lar{\leftarrow}
\newcommand\Lar{\Leftarrow}
\newcommand\xlar{\xleftarrow}
\newcommand\lrar{\leftrightarrow}
\newcommand\Lrar{\Leftrightarrow}
\newcommand\ideal{\trianglelefteq}
\newcommand\ai{\text{ a.\^{i}. }}
\newcommand\sk{(\textit{Schi\c{t}\u{a}}) }
\newcommand\rhup{\rightharpoonup}
\newcommand\rhdn{\rightharpoondown}
\newcommand\lhup{\leftharpoonup}
\newcommand\lhdn{\leftharpoondown}
\newcommand\vid{\emptyset}
\newcommand\wh{\widehat}
\newcommand\ol{\overline}
\newcommand\NN{\mathbb{N}}
\newcommand\RR{\mathbb{R}}
\newcommand\ZZ{\mathbb{Z}}
\newcommand\QQ{\mathbb{Q}}
\newcommand\CC{\mathbb{C}}

%% STILURILE MELE
\newtheoremstyle{dotless}{}{}{\itshape}{}{\bfseries}{:}{ }{}

\theoremstyle{dotless}
\newtheorem{theorem}{Teorem\u{a}}[section]
\newtheorem{proposition}{Propozi\c{t}ie}[section]
\newtheorem{lemma}{Lem\u{a}}[section]
\newtheorem{corollary}{Corolar}[section]
\newtheorem{exercise}{Exerci\c{t}iu}

\newtheoremstyle{dotlessdr}{}{}{}{}{\bfseries}{:}{ }{}
\theoremstyle{dotlessdr}
\newtheorem{example}{Exemplu}[section]
\newtheorem{definition}{Defini\c{t}ie}[section]
\newtheorem{remark}{Observa\c{t}ie}[section]




\begin{document}

\begin{center}
  {\large\textbf{Seminar 7 \\
      Integrale de suprafață. Formule integrale}}
\end{center}

\vspace{1cm}

\section{Integrale de suprafață}
\label{sec:teorie}

Integralele de suprafață sînt analogul celor curbilinii, dar, în locul unei curbe
în lungul căreia se face integrarea, ne vom baza pe o suprafață din $\RR^3$.

Tot ca în cazul integralelor curbilinii, integralele de suprafață sînt de două tipuri.

\textbf{Integrala de suprafață de speța întîi} are forma și se calculează după formula:
\[
  \iint_\Sigma \mathcal{F}(x, y, z) d\sigma =
  \iint_D \mathcal{F}(x(u,v), y(u,v), z(u,v)) \sqrt{EG - F^2} dudv,
\]
unde $\Sigma$ este suprafața pe care integrăm, care devine domeniul $D$ după aplicarea
formulei, cu elementul de arie $d \sigma$, care devine $dudv$, fiecare dintre cele trei
coordonate $x, y, z$ fiind parametrizate după $u$ și $v$, iar coeficienții $E, F, G$
se calculează astfel:

\begin{align*}
  E &= \Big( \frac{\partial x}{\partial u} \Big)^2 + \Big(\frac{\partial y}{\partial u}\Big)^2 +
      \Big(\frac{\partial z}{\partial u}\Big)^2 \\
  F &= \frac{\partial x}{\partial u} \cdot \frac{\partial x}{\partial v} +
      \frac{\partial y}{\partial u} \cdot \frac{\partial y}{\partial v} +
      \frac{\partial z}{\partial u} \cdot \frac{\partial z}{\partial v} \\
  G &= \Big(\frac{\partial x}{\partial v} \Big)^2 + \Big(\frac{\partial y}{\partial v}\Big)^2 +
      \Big(\frac{\partial z}{\partial v} \Big)^2.
\end{align*}

O definiție alternativă este următoarea. Fie $\Phi : D \to \RR^3$ o pînză parametrizată
și fie $\Sigma = \Phi(D)$ imaginea ei, iar $F : U \to \RR$ o funcție continuă pe
imaginea pînzei. Integrala de suprafață a lui $F$ pe $\Sigma$ este, prin definiție:
\[
  \int_\Sigma F(x, y, z) d \sigma =
  \iint_D F(\Phi(u, v)) \cdot \Big|\Big| \frac{\partial \Phi}{\partial u} \times %
  \frac{\partial \Phi}{\partial v} \Big|\Big| dudv.
\]

Într-un caz particular, în care suprafața regulată $\Sigma$ este dată sub o formă
explicită $z = f(x, y)$, cu $f : D \seq \RR^2 \to \RR$, atunci formula se simplifică
și devine:
\[
  \iint_\Sigma \mathcal{F}(x, y, z) d \sigma =
  \iint_D \mathcal{F}(x, y, f(x, y)) \sqrt{1 + p^2 + q^2} dx dy,
\]
unde coeficienții au forma simplificată:
\[
  p = \frac{\partial z}{\partial x}, \quad q = \frac{\partial z}{\partial y}.
\]

Alte cazuri particulare sînt:
\begin{enumerate}[(a)]
\item Pentru $F = 1$, obținem \emph{aria suprafeței $\Sigma$};
\item Dacă $F \geq 0$ este densitatea unei plăci $\Sigma$, atunci \emph{masa ei}
  se calculează prin $M = \int_\Sigma F d \sigma$, iar \emph{coordonatele centrului %
  de greutate}: $x_i^G = \frac{1}{M} \int_\Sigma x_i F d \sigma$.
\end{enumerate}

Într-o \textbf{interpretare vectorială}, putem considera $\vec{F}$ un cîmp vectorial,
atunci \emph{fluxul} său prin suprafața $\Sigma$ este dat de o integrală de suprafață
de prima speță, cu forma:
\[
  \iint_\Sigma \vec{F} \cdot \vec{n} d \sigma,
\]
unde $\vec{n}$ este vectorul normal la suprafață, calculat prin
$\vec{n} = \frac{\nabla F}{||\nabla F||}$, iar $d \sigma$ este elementul de suprafață.

Dacă avem cîmpul vectorial dat parametric, atunci formula pentru flux devine:
\[
  \Phi_\Sigma (\vec{F}) = \iint_D F(x(u, v), y(u, v), z(u, v)) \cdot
  (\vec{T}_u \times \vec{T}_v) dS,
\]
unde $\vec{T}_u, \vec{T}_v$ sînt vectorii tangenți după cele două direcții, calculați
ca derivatele parțiale $F_u$ și $F_v$, iar $dS$ este elementul de arie, adică $dudv$.

\vspace{1cm}

Pentru \textbf{integralele de suprafață de speța a doua}, considerăm o 2-formă
diferențială:
\[
  \omega = P dy \wedge dz + Q dz \wedge dx + R dx \wedge dy
\]
și luăm o pînză parametrizată:
\[
  \Phi : D \to \RR^3, \quad \Phi(u, v) = (X(u, v), Y(u, v), Z(u, v)).
\]

Integrala pe suprafața orientată $\Sigma$ a formei diferențiale $\omega$ este
definită prin:
\[
  \int_\Sigma \omega = \iint_D \Big( (P \circ \Phi) \frac{D(Y,Z)}{D(u, v)} +
  (Q \circ \Phi) \frac{D(Z, X)}{D(u,v)} + (R \circ \Phi) \frac{D(X, Y)}{D(u,v)}\Big)dudv,
\]
unde $\frac{D(Y,Z)}{D(u,v)}$ etc. sînt jacobienii funcțiilor $X, Y, Z$ în raport cu
variabilele $u, v$.

Într-o formă simplificată, dacă privim $\vec{n} = \frac{\nabla \Phi}{||\nabla \Phi||}$,
iar $\Phi = (X, Y, Z)$, atunci putem scrie:
\[
  \iint_\Sigma \omega = \iint_D
  \begin{vmatrix}
    P \circ \Phi & Q \circ \Phi & R \circ \Phi \\
    X_u & Y_u & Z_u \\
    X_v & Y_v & Z_v
  \end{vmatrix} dudv,
\]
unde am folosit notația simplificată pentru derivate parțiale, i.e.\
$X_u = \frac{\partial X}{\partial u}$ etc.

%%%%%%%%%%%%%%%%%%%%%%%%%%%%%%%%%%%%%%%%%%%%%%%%%%%%%%%%%%%%%%%%%%%%%%

\section{Exerciții}
\label{sec:exc}

1. În fiecare din exemplele următoare, considerăm:
\[
  D \ni (u,v) \mapsto \Phi(u,v) = (X(u,v), Y(u,v), Z(u,v)) \in \RR^3
\]
o pînză parametrizată. Să se calculeze vectorii tangenți la suprafață și
versorul normalei la suprafață.

\begin{enumerate}[(a)]
\item \emph{Sfera}: Fie $R > 0$ și $\Phi : [0, \pi] \times [0, 2\pi) \to \RR^3$:
  \[
    \Phi(\theta, \varphi) = (R\sin \theta \cos \varphi, R \sin \theta \sin \varphi, R \cos \theta);
  \]
\item \emph{Paraboloidul}: Fie $a > 0, h > 0$ și $\Phi : [0, h] \times [0, 2\pi) \to \RR^3$:
  \[
    \Phi(u, v) = (au\cos v, au\sin v u^2);
  \]
\item \emph{Elipsoidul}: Fie $a, b, c > 0$ și $\Phi : [0, \pi] \times [0, 2\pi) \to \RR^3$:
  \[
    \Phi(\theta, \varphi) = (a \sin\theta \cos \varphi, b \sin \theta \sin \varphi, c \cos \theta);
  \]
\item \emph{Conul}: Fie $h > 0, \Phi : [0, 2\pi) \times [0, h] \to \RR^3$:
  \[
    \Phi(u, v) = (v \cos u, v \sin u, v);
  \]
\item \emph{Cilindrul}: Fie $a > 0, 0 \leq h_1 \leq h_2, \Phi : [0, 2\pi) \times [h_1, h_2] \to \RR^3$:
  \[
    \Phi(\varphi, z) = (a \cos \varphi, a \sin \varphi, z);
  \]
\item \emph{Torul}: Fie $0 < a < b, \Phi [0, 2\pi) \times [0, 2\pi) \to \RR^3$:
  \[
    \Phi(u, v) = ((a + b\cos u) \cos v, (a + b \cos u) \sin v, b \sin u).
  \]
\end{enumerate}

\emph{Indicație:} Vectorii tangenți la suprafață sînt $\frac{\partial \Phi}{\partial u}$ și
$\frac{\partial \Phi}{\partial v}$, iar versorul normalei este:
\[
  \frac{1}{\Big|\Big| \frac{\partial \Phi}{\partial u} \times \frac{\partial \Phi}{\partial v}\Big|\Big|}
    \cdot \frac{\partial \Phi}{\partial u} \times \frac{\partial \Phi}{\partial v}.
\]

\vspace{1cm}

2. Fie $\omega = Pdy \wedge dz + Qdz \wedge dx + Rdx \wedge dy$ o 2-formă diferențială,
iar $\Sigma$, imaginea unei pînze parametrizate. Să se calculeze integrala de
suprafață $\int_\Sigma \omega$ în cazurile:
\begin{enumerate}[(a)]
\item $\omega = ydy \wedge dz + zdz \wedge dx + xdx \wedge dy$, unde:
  \[
    \Sigma : \begin{cases}
      X(u, v) &= u \cos v \\
      Y(u, v) &= u \sin v \\
      Z(u, v) &= cv
    \end{cases}
  \]
  cu domeniul $(u, v) \in [a, b] \times [0, 2\pi)$;
\item $\omega = xdy \wedge dz + ydz \wedge dx + zdx \wedge dy$, unde:
  \[
    \Sigma : x^2 + y^2 + z^2 = R^2;
  \]
\item $\omega = yzdy \wedge dz + zxdz \wedge dx + xydx \wedge dy$, unde:
  \[
    \Sigma : \frac{x^2}{a^2} + \frac{y^2}{b^2} + \frac{z^2}{c^2} = 1;
  \]
\item $\omega = xdy \wedge dz + ydz \wedge dx$, unde:
  \[
    \Sigma: x^2 + y^2 = z^2, \quad z \in [1,2];
  \]
\item $\omega = (y + z) dy \wedge dz + (x + y) dx \wedge dy$, unde:
  \[
    \Sigma: x^2 + y^2 = a^2, \quad a \in [0, 1]
  \]
\end{enumerate}

\vspace{1cm}

3. Să se calculeze integrala de suprafață de prima speță:
\[
  \iint_\Sigma F(x, y, z) d\sigma
\]
în următoarele cazuri:
\begin{enumerate}[(a)]
\item $F(x, y, z) = |xyz|$, iar $\Sigma: x^2 + y^2 = z^2, z \in [0, 1]$;
\item $F(x, y, z) = y\sqrt{z}$, iar $\Sigma: x^2 + y^2 = 6z, z \in [0, 2]$;
\item $F(x, y, z) = z^2$, iar $\Sigma = \{(x, y, z) \mid z = \sqrt{x^2 + y^2}, %
  x^2 + y^2 - 6y \leq 0 \}$.
\end{enumerate}

\vspace{1cm}

4. Folosind integralele de suprafață, să se calculeze ariile suprafețelor:
\begin{enumerate}[(a)]
\item sfera de rază $R$;
\item conul $z^2 = x^2 + y^2, z \in [0, h]$;
\item paraboloidul $z = x^2 + y^2, z \in [0, h]$.
\end{enumerate}

\vspace{1cm}

5. Să se calculeze aria suprafeței $\Sigma$ în următoarele cazuri:
\begin{enumerate}[(a)]
\item $\Sigma : 2z = 4 - x^2 - y^2, z \in [0, 1]$;
\item $\Sigma$ este submulțimea de pe sfera $x^2 + y^2 + z^2 = 1$, situată în interiorul
  conului $x^2 + y^2 = z^2$;
\item $\Sigma$ este submulțimea de pe sfera $x^2 + y^2 + z^2 = R^2$, situată în interiorul
  cilindrului $x^2 + y^2 - Ry = 0$;
\item $\Sigma$ este submulțimea de pe paraboloidul $z = x^2 + y^2$, situată în interiorul
  cilindrului $x^2 + y^2 = 2y$.
\end{enumerate}

\vspace{1cm}

6. Să se calculeze fluxul cîmpului vectorial $\vec{V}$ prin suprafața $\Sigma$
în următoarele cazuri:
\begin{enumerate}[(a)]
\item $\vec{V} = x\vec{i} + y\vec{j} + z \vec{k}$, iar $\Sigma: z^2 = x^2 + y^2$,
  cu $z \in [0, 1]$;
\item $\vec{V} = y\vec{i} + x \vec{j} + z^2 \vec{k}$, iar $\Sigma: z = x^2 + y^2$, cu
  $z \in [0, 1]$;
\item $\vec{V} = \frac{1}{\sqrt{x^2 + y^2}} (y\vec{i} - x \vec{j} + \vec{k})$, iar
  $\Sigma: z = 4 - x^2 - y^2, z \in [0, 1]$.
\end{enumerate}


%%%%%%%%%%%%%%%%%%%%%%%%%%%%%%%%%%%%%%%%%%%%%%%%%%%%%%%%%%%%%%%%%%%%%%

\section{Formula Green-Riemann}
\label{sec:g-r}

Fie $(K, \partial K)$ un compact cu bord orientat inclus în $\RR^2$ și considerăm
o 1-formă diferențială de clasă $\kal{C}^1$ pe o vecinătate a lui $K$. Atunci
are loc \emph{formula Green-Riemann}:
\[
  \int_{\partial K} Pdx + Qdy = \iint_K \Big(\frac{\partial Q}{\partial x} - %
  \frac{\partial P}{\partial y} \Big) dx dy.
\]

O consecință imediată este o formulă pentru arie:
\[
  \mathcal{A}(K) = \frac{1}{2} \int_{\partial K} xdy - ydx.
\]

%%%%%%%%%%%%%%%%%%%%%%%%%%%%%%%%%%%%%%%%%%%%%%%%%%%%%%%%%%%%%%%%%%%%%%

\section{Formula Gauss-Ostrogradski}
\label{sec:g-o}

Considerăm $K$ o mulțime compactă, cu bord orientat după normala exterioară. Atunci,
pentru orice 2-formă diferențială:
\[
  \omega = Pdy \wedge dz + Q dz \wedge dx + R dx \wedge dy
\]
de clasă $\kal{C}^1$ pe o vecinătate a lui $K$ are loc egalitatea:
\[
  \int_{\partial K} P dy \wedge dz + Q dz \wedge dx + R dx \wedge dy = %
  \iiint_K \frac{\partial P}{\partial x} + \frac{\partial Q}{\partial y} + %
  \frac{\partial R}{\partial z} dxdydz.
\]

În notație vectorială, dacă $\vec{V} = P \vec{i} + Q \vec{j} + R \vec{k}$ este
cîmpul vectorial asociat 2-formei diferențiale $\omega$, atunci formula de mai sus
poate fi scrisă:
\[
  \int_{\partial K} \vec{V} \cdot \vec{n} d \sigma = \iiint_K \mathrm{div}\vec{V} dxdydz,
\]
unde $\vec{n}$ este normala exterioară la $\partial K$. În membrul stîng avem
fluxul cîmpului $\vec{V}$ prin suprafața $\partial K$, motiv pentru care formula
Gauss-Ostrogradski se mai numește \emph{formula flux-divergență}.

%%%%%%%%%%%%%%%%%%%%%%%%%%%%%%%%%%%%%%%%%%%%%%%%%%%%%%%%%%%%%%%%%%%%%%

\section{Formula lui Stokes}
\label{sec:stokes}

Fie $\Sigma$ o suprafață cu bord, orientată și fie:
\[
  \alpha = Pdx + Qdy + Rdz
\]
o 1-formă diferențială de clasă $\kal{C}^1$ pe o vecinătate a lui $\Sigma$.
Atunci are loc formula:
\begin{align*}
  \int_{\partial \Sigma} Pdx + Qdy + Rdz &= \int_\sigma \Big( \frac{\partial R}{\partial y} -
                                           \frac{\partial Q}{\partial z} \Big) dy \wedge dz +
                                           \Big( \frac{\partial P}{\partial z} -
                                           \frac{\partial R}{\partial x} \Big) dz \wedge dx +
                                           \Big( \frac{\partial Q}{\partial x} -
                                           \frac{\partial P}{\partial y}\Big) dx \wedge dy.
\end{align*}

În notație vectorială, dacă $\vec{V} = P \vec{i} + Q \vec{j} + R \vec{k}$ este cîmpul
vectorial asociat formei diferențiale $\alpha$, atunci formula lui Stokes se scrie:
\[
  \int_{\partial \Sigma} \vec{V} \cdot d\vec{r} = \int_\Sigma (\nabla \times \vec{V}) %
  \cdot \vec{n} d \sigma,
\]
unde $\nabla \times \vec{V} = \mathrm{rot}\vec{V}$ se numește \emph{rotorul} cîmpului
vectorial $\vec{V}$, calculat cu ajutorul produsului vectorial.

%%%%%%%%%%%%%%%%%%%%%%%%%%%%%%%%%%%%%%%%%%%%%%%%%%%%%%%%%%%%%%%%%%%%%%

\section{Exerciții}
\label{sec:exc-f}

1. Să se calculeze direct și folosind formula Green-Riemann integrala curbilinie
$\int_\Gamma \alpha$ în următoarele cazuri:
\begin{enumerate}[(a)]
\item $\alpha = y^2 dx + xdy$, unde $\Gamma$ este pătratul cu vîrfurile
  $A(0, 0), B(2, 0), C(2,2), D(0, 2)$;
\item $\alpha = ydx + x^2 dy$, unde $\Gamma$ este cercul cu centrul în origine și
  rază 2;
\item $\alpha = ydx - xdy$, unde $\Gamma$ este elipsa de semiaxe $a$ și $b$ și de
  centru $O$.
\end{enumerate}

\vspace{1cm}

2. Să se calculeze integralele, direct, apoi aplicînd formula Green-Riemann:
\begin{enumerate}[(a)]
\item $\ds\int_\Gamma e^{\frac{x^2}{a^2} + \frac{y^2}{b^2}} (-ydx + xdy)$, unde
  $\Gamma = \{(x, y) \mid \frac{x^2}{a^2} + \frac{y^2}{b^2} = 1 \}$;
\item $\ds\int_\Gamma xy dx + \frac{x^2}{2} dy$, unde $\Gamma$ este obținută prin:
  \[
    \Gamma = \{x^2 + y^2 = 1, x \leq 0 \leq y\} \cup \{x + y = -1, x,y \leq 0 \}.
  \]
\end{enumerate}

\vspace{1cm}

3. Să se calculeze circulația cîmpului vectorial $\vec{V}$ pe curba $\Gamma$ în
cazurile:
\begin{enumerate}[(a)]
\item $\vec{V} = y^2 \vec{i} + xy\vec{j}$, unde:
  \[
    \Gamma = \{x^2 + y^2 = 1, y > 0 \} \cup \{ y = x^2 - 1, y \leq 0 \};
  \]
\item $\vec{V} = e^x \cos y \vec{i} - e^x \sin y \vec{j}$, unde $\Gamma$ este
  o curbă arbitrară din semiplanul superior, ce unește punctele $A(1, 0), B(-1, 0)$,
  cu sensul de la $A$ către $B$.
\end{enumerate}

\vspace{1cm}

4. Să se calculeze integrala de suprafață $\int_\Sigma \omega$ în următoarele cazuri:
\begin{enumerate}[(a)]
\item $\omega = x^2 dy \wedge dz - 2xy dz \wedge dx + z^3 dx \wedge dy$, iar mulțimea
  $\Sigma = \{x^2 + y^2 + z^2 = 9 \}$;
\item $\omega = yz dy \wedge dz - (x + z) dz \wedge dx + (x^2 + y^2 + 3z) dx \wedge dy$,
  iar mulțimea:
  \[
    \Sigma = \{ x^2 + y^2 = 4 - 2z, z \geq 1 \} \cup \{x^2 + y^2 \leq 4 - 2z, z = 1 \};
  \]
\item $\omega = x(z + 3) dy \wedge dz + yz dz \wedge dx - (z + z^2) dx \wedge dy$,
  iar mulțimea:
  \[
    \Sigma = \{ x^2 + y^2 + z^2 = 1, z \geq 0\}.
  \]
\end{enumerate}

\vspace{1cm}

5. Să se calculeze, folosind formula lui Stokes, integrala curbilinie
$\int_\Gamma \alpha$, în următoarele cazuri:
\begin{enumerate}[(a)]
\item $\alpha = (y - z) dx + (z - x)dy + (x - y) dz$, iar mulțimea $\Gamma: z = x^2 + y^2, z = 1$;
\item $\alpha = ydx + zdy + xdz$, iar mulțimea:
  \[
    \Gamma: x^2 + y^2 + z^2 = 1, \quad x + y + z = 0.
  \]
\end{enumerate}

\vspace{1cm}

6. Să se calculeze circulația cîmpului vectorial:
\[
  \vec{V} = (y^2 + z^2)\vec{i} + (x^2 + z^2)\vec{j} + (x^2 + y^2)\vec{k},
\]
pe curba $\Gamma: x^2 + y^2 + z^2 = R^2, ax + by + cz = 0$.


\end{document}
